\part{Discovery and Synthesis}

% =====================================================================
\chapter{arXiv Browse}
\label{ch:browse}

The \cmd{/browse} route is a page designed for researchers who like to "browse arXiv." It organizes
papers by the official arXiv category tree, letting you explore the latest papers in a field without
having specific keywords in mind.

\section{Category Tree}

\litfigure{07-browse-en.png}{arXiv Browse. The left sidebar shows the category tree (top-level groups such as cs / math / physics);
the main area displays the latest papers in the selected category, with pagination and import support.}

LitFolio includes a built-in Chinese--English mapping for all 150+ arXiv categories. Clicking any
leaf node (e.g.\ \cmd{cs.LG} for Machine Learning, \cmd{physics.optics} for Optics) triggers an
arXiv API request to fetch the 50 most recent papers.

\section{Pagination and "Exhausted"}

Scrolling to the bottom and clicking "Load More" fetches the next page (50 items per page). When a
category has low publication volume and two consecutive requests return no new entries, LitFolio
displays "Exhausted" and stops requesting---this avoids unnecessary API calls.

\section{Importing}

Clicking any paper opens a detail drawer showing the full abstract, authors, submission date, and
original arXiv link. The "Import" button in the bottom-right corner of the drawer follows the same
arXiv ID import workflow described in Chapter 3---the only difference is that the ID is already
filled in.

\begin{tipbox}
The arXiv Browse page itself does \textbf{not} save the list to the local database---it queries
the live API every time. If you want to follow a category's new papers long-term, consider using
RSS subscriptions (see the next chapter), which work offline, support translation, and track read
status.
\end{tipbox}

% =====================================================================
\chapter{RSS Feeds}
\label{ch:feeds}

The \cmd{/feeds} route lets you track arXiv sub-categories, journal articles, and research group
homepages via RSS. All fetched entries are stored in SQLite, with unread counts and read status
saved locally.

\section{Feed Management}

\litfigure{08-feeds-en.png}{Feed list. Left: sources (arXiv cs.LG / Nature Photonics / Optica, etc.); Right: latest entries
for the selected source, with unread badges, translation buttons, and import buttons.}

To add a new feed source, click the "+" button and enter the RSS URL and a display name. LitFolio
imposes no restrictions on feed sources---any valid RSS/Atom feed works. See Appendix C for
commonly used sources.

\section{Conditional GET Refresh}

LitFolio uses conditional GET (\cmd{If-Modified-Since} / \cmd{ETag}) to fetch feeds. When the
remote has not been updated, the server returns 304 Not Modified---this is why clicking "Refresh"
sometimes instantly responds with "Already up to date." In practice:

\begin{itemize}
  \item Refreshing frequently is \textbf{cheap}---it will not trigger rate limiting on the server.
  \item Genuine "nothing new" and "network/server error" states display the same message; you need
    to check the error badge to tell them apart.
\end{itemize}

\section{Entry Details and Translation}

\litfigure{09-feeds-detail-en.png}{After selecting a feed item, the right drawer shows the full abstract; you can translate
with one click, mark as read, or import directly (equivalent to pasting an arXiv ID on the Import page).}

The "Translate" button works the same way as in the library: it invokes the \cmd{translate} task
(bound to the translation model you configured). Results are cached in the SQLite \cmd{feed\_items}
table and load instantly on subsequent views.

\begin{warnbox}
RSS translation consumes tokens for every entry. If your subscribed sources update frequently
(e.g., multiple arXiv categories), consider binding the translation task to a \textbf{local Ollama
open-source chat model} or a \textbf{low-cost domestic API} to avoid runaway costs.
\end{warnbox}

\section{Direct Import from Feed Items}

The "Import" button in the detail drawer follows the exact same workflow as the arXiv ID import
process in Chapter 3---fetching full metadata and attempting to download the PDF. After import, the
entry remains in the feed list but is marked as "Imported."

% =====================================================================
\chapter{Topic Discovery}
\label{ch:topic}

The \cmd{/topic} route is one of LitFolio's advanced capabilities, split into two tabs:
\textbf{Search Recall} and \textbf{Survey Generation}. They solve different problems.

\section{The Two Tabs}

\begin{description}
  \item[Search Recall] You already have a rough research direction and want to find all relevant
    papers from your existing library. Workflow: query expansion $\to$ multi-path FTS5 retrieval
    $\to$ rank by citation count. Fully offline, no external network calls.
  \item[Survey Generation] You want a "bird's-eye view" of a new direction, but your library may
    not contain any papers on that topic. Workflow: LLM decomposes into sub-topics $\to$ Semantic
    Scholar fetches real papers $\to$ LLM annotates must-reads.
\end{description}

\section{Search Recall}

\litfigure{10-topic-search-en.png}{Search Recall tab. After entering a query, LitFolio expands it into several related phrases,
retrieves them via multi-path FTS5, and ranks results by citation count. Clicking a result opens a drawer with the TL;DR.}

\subsection{Query Expansion}

Entering "retrieval augmented generation" gets expanded to queries like "RAG / retrieval-augmented
/ vector retrieval / chunk + retrieval / Lewis 2020..." and so on---a dozen or so related queries.
This step uses the LLM configured for the \cmd{search} task, with the goal of recalling papers that
use different phrasing to describe the same concept.

\subsection{Query Expansion Algorithm Details}

Expansion is handled by the \cmd{query\_expand} module. It does not pass the user's raw input
directly to S2; instead, it first asks the LLM to translate the (possibly Chinese) broad topic
into 2--4 precise English search terms. The core constraints are written in the system prompt:

\begin{itemize}
  \item Output only 2--4 search terms---fewer is better than more, as multi-term OR queries only
    inflate noise.
  \item Choose terms that domain experts would actually type, rather than piling on keywords from
    adjacent sub-fields.
  \item If two candidate terms would recall the same set of papers (e.g., "attosecond pulses" and
    "attosecond laser"), keep only the more standard one.
  \item Chinese input is translated to English before searching.
  \item Remove vague modifiers (research / study / novel / recent, etc.).
\end{itemize}

The LLM returns JSON: \cmd{\{"terms": [...]\}}. The parser has three-level fault tolerance: it
first attempts direct JSON parsing; on failure, it extracts the first \cmd{\{...\}} substring
(handling markdown fences and conversational wrappers); as a last resort, it splits by comma or
newline. At most 4 terms are used.

\subsection{Time Window and Sorting}

Results are sorted by citation count by default; the dropdown in the top-right corner allows
switching to "Newest," "Relevance," or "Must-Read First." The time window (last 1/3/5 years)
filters out older papers.

\section{Survey Generation}

\litfigure{11-topic-survey-en.png}{Survey Generation tab. Enter an open topic, and LitFolio uses an LLM to decompose it into
3 sub-fields, fetches 4 representative papers per sub-field from Semantic Scholar, then has the LLM mark "Must-Read" with a one-line justification.}

\subsection{Pipeline}

Survey generation is a four-step pipeline, with each step's progress pushed to the frontend in
real time via IPC events.

\subsubsection{Why Not Let the LLM List Papers Directly}

Asking an LLM to list papers directly leads to fabricated DOIs, authors, and
years---hallucination rates are extremely high. LitFolio separates "planning"
from "retrieval": the LLM only handles decomposing sub-fields and providing
search terms. Actual paper retrieval from Semantic Scholar is done by program code, and the
LLM then annotates the real papers.

\begin{enumerate}
  \item \textbf{Planning}---Decompose the open topic into 3 sub-fields (e.g., "extreme ultrashort
    pulse laser" might decompose into chirped pulse amplification / harmonic generation / optical
    parametric processes).
  \item \textbf{Grounding}---For each sub-field, call the Semantic Scholar Graph API to fetch 4
    highly cited papers. These are real references, not LLM fabrications.
  \item \textbf{Annotating}---The LLM marks each paper as "Must-Read" with a one-line
    justification (explaining why this paper is important for the sub-field).
  \item \textbf{Done}---Save the complete survey to the SQLite \cmd{topic\_surveys} table for
    later retrieval.
\end{enumerate}

\subsubsection{Phase 1: LLM Planning}

The LLM receives a (possibly Chinese) research topic and outputs a JSON structure: 4--7 sub-fields,
each containing a name, active year range, a Chinese narrative, 2--4 English S2 search terms, and
an optional PI name hint. Additionally, 5--10 cross-sub-field key\_PIs (important researchers) are
included. Chinese input is translated into English search terms within the prompt.

\subsubsection{Phase 2: Semantic Scholar Grounding}

For each sub-field, the program executes the following steps:

\begin{enumerate}
  \item \textbf{Term-based retrieval}---Each search\_term triggers one call to the S2
    \cmd{bulk\_by\_citations} API, returning papers sorted by citation count with year filtering.
    Each term fetches \cmd{topk $\times$ 3} (at least 15) results to ensure sufficient candidates
    after deduplication.
  \item \textbf{PI supplementary retrieval}---Combine PI names with search\_terms to query S2;
    each PI is combined with at most 2 search terms.
  \item \textbf{Deduplication}---Within each sub-field, deduplicate by paperId (prefer paperId;
    use DOI if paperId is unavailable). When duplicates exist, keep the version with higher citation
    count.
  \item \textbf{Relevance filtering}---The paper title and abstract must contain at least 2
    stop-word-removed search term tokens, or directly contain a complete search phrase. S2
    occasionally returns irrelevant papers; this filtering layer specifically catches them.
  \item \textbf{DOI filtering + sorting}---Papers without a DOI are discarded (to ensure
    traceability). The remaining papers are sorted by citation count in descending order and the
    top-K are selected.
\end{enumerate}

Cross-sub-field deduplication uses a "first come, first served" strategy: if a paper matches in
multiple sub-fields, it is kept only in the sub-field processed first.

\subsubsection{S2 API Rate Control}

LitFolio has a built-in token bucket rate limiter: at most 80 requests per 300-second window (S2's
unauthenticated limit is 100 requests per 5 minutes; LitFolio leaves a 20\% margin). It is
implemented with a global atomic counter that resets automatically when the window expires. When the
limit is exceeded, a clear error message "try again later" is returned instead of silent failure.
During survey generation, multiple search\_terms within a sub-field are called serially. If there
are many sub-fields and search terms, rate limiting may be triggered---this is why survey generation
sometimes takes 30--60 seconds.

\subsubsection{Phase 3: LLM Annotation (Optional)}

The real paper list from all sub-fields (id + title + year + abstract) is sent to the LLM with the
following requirements:

\begin{itemize}
  \item Write 1--2 sentences in Chinese explaining "why this matters" for each paper.
  \item Mark 8--12 papers across all sub-fields as "Must-Read," ensuring coverage across different
    sub-fields (do not cluster them in a single sub-field).
\end{itemize}

Phase 3 is optional---if this LLM call fails, the survey is still usable; it simply lacks
annotations and must-read markers.

\subsection{Saving and Restoring}

Each generation is automatically saved. The "History" button lets you return to a previous survey
without re-calling the LLM or APIs.

\begin{tipbox}
Chinese input is fully supported. LitFolio translates the Chinese topic into English search terms
during the planning step before querying S2, so "extreme ultrashort pulse laser" becomes something
like "ultrashort intense laser pulses" for literature retrieval. The entire process typically takes
30--60 seconds.
\end{tipbox}

\section{Topic Alerts}

Topic Alerts are an automated extension of Topic Discovery---once you set a research direction,
LitFolio periodically monitors for newly published relevant papers. Unlike RSS subscriptions, which
track \textbf{sources} (an arXiv category or journal), Topic Alerts track \textbf{topics}
(cross-source research directions).

\litfigure{27-topic-alerts-en.png}{Topic Alert management interface. Set the query, frequency, and target folder;
the system runs automatically and reports newly discovered papers.}

\subsection{Creating an Alert}

Creating an alert requires the following settings:

\begin{description}
  \item[Query] A natural language description of the research topic (same input format as Topic
    Discovery), e.g., "chirped pulse amplification for extreme ultrashort pulse lasers."
  \item[Frequency] Daily / Weekly / On startup.
  \item[Target Folder] Optional---papers imported automatically will be placed in the specified
    folder.
  \item[Auto Import] When enabled, newly discovered papers are automatically imported into the
    library.
\end{description}

\subsection{Alert Execution Flow}

When an alert triggers, LitFolio executes the following flow:

\begin{enumerate}
  \item \textbf{LLM Planning}---Translate the query into 2--4 English search terms (using the same
    logic as Topic Discovery's \cmd{query\_expand}).
  \item \textbf{Semantic Scholar Retrieval}---Call the S2 API for each search term, fetching the
    latest papers sorted by citation count.
  \item \textbf{Deduplication Filtering}---Compare the retrieval results against two lists:
    \begin{itemize}
      \item Already-imported papers (DOI / arXiv ID match) $\to$ excluded.
      \item Papers shown in previous alerts (\cmd{topic\_alert\_results} table) $\to$ excluded.
    \end{itemize}
    Only papers that are \textbf{neither imported nor previously shown} are kept.
  \item \textbf{Store Results}---New papers are written to the \cmd{topic\_alert\_results} table
    with \cmd{seen = 0}.
  \item \textbf{Notification}---The topic entry in the navigation bar displays an unread count
    badge.
\end{enumerate}

\subsection{Auto Import}

If "Auto Import" is enabled, the following additional steps execute after step 4:

\begin{enumerate}
  \setcounter{enumi}{4}
  \item \textbf{Batch Import}---For each new paper, invoke the arXiv/DOI import workflow to fetch
    metadata and PDF.
  \item \textbf{Place in Target Folder}---If a target folder is specified, imported papers are
    automatically placed into that folder.
\end{enumerate}

\subsection{Management and Viewing}

Manage all alerts in Settings $\to$ Topic Alerts:

\begin{itemize}
  \item View the history of results for each alert.
  \item Mark results as read (\cmd{seen = 1}).
  \item Edit an alert's query, frequency, or target folder.
  \item Pause or delete an alert.
\end{itemize}

\begin{warnbox}
Auto Import consumes LLM tokens (for metadata translation, etc.). If you have multiple
high-frequency alerts with Auto Import enabled, consider binding the related tasks to a low-cost
model.
\end{warnbox}

% =====================================================================
\chapter{Library Q\&A (RAG)}
\label{ch:ask}

The \cmd{/ask} route turns LitFolio into \textbf{your personal knowledge base} for asking
questions. Unlike generic ChatGPT, answers are based solely on papers you have already imported,
and every factual claim is accompanied by a \cmd{[N]} citation reference.

\section{How It Works}

\litfigure{12-ask-empty-en.png}{Ask page initial state. Three cards explain the workflow: 1) question rewriting;
2) multi-path FTS5 retrieval; 3) LLM answer with citations.}

The entire process has four steps:

\begin{enumerate}
  \item \textbf{LLM Question Rewriting}---This step is the key to the entire RAG pipeline. Throwing
    a Chinese question directly into FTS5's AND-of-tokens path yields essentially zero recall
    (FTS5 \cmd{unicode61} tokenizes Chinese character by character, so "neural network" becomes
    "neural AND network" in individual characters). Therefore, the \cmd{query\_expand} module first
    rewrites the natural language question into 2--4 English search terms. If rewriting fails
    (offline, invalid key, model not configured), it does not block the flow---it degrades to
    searching with the original question directly.
  \item \textbf{Multi-Path FTS5 Retrieval + Merged Scoring}---Each search term runs an independent
    FTS5 BM25 search (each term over-fetches \cmd{limit $\times$ 3}, at least 8 results), then
    results are merged by paper\_id. The sorting rule during merging is: \textbf{number of matched
    terms} (papers matched by more search terms rank higher) $>$ descending year $>$ descending
    import time. This way, papers that are "relevant from multiple semantic angles" take priority
    over papers that score high on only one term. If all expanded queries return empty, a
    \textbf{degraded search chain} kicks in: raw question AND search $\rightarrow$ all-term OR
    fuzzy matching $\rightarrow$ CJK keyword split fuzzy search $\rightarrow$ final fallback to
    recent library papers, ensuring the LLM always has some context to work with.
  \item \textbf{Evidence Assembly}---Each recalled paper is compressed into a snippet of no more
    than 1800 characters, assembled by priority: TL;DR $\to$ Abstract $\to$ Problem $\to$ Method
    $\to$ Comparison $\to$ Limitations $\to$ Key Findings $\to$ User highlights (up to 3, each 240
    characters). All snippets combined do not exceed 14,000 characters; excess is truncated.
  \item \textbf{LLM Answer}---The system prompt enforces: answer only based on the provided
    evidence, annotate each fact with a \cmd{[N]} citation reference; when the sources are
    insufficient, the LLM says so explicitly and describes what kind of paper would be needed
    rather than guessing. Answers are rendered in rich Markdown format (bold, italic, lists,
    headings).
\end{enumerate}

\subsection{@-Mention Pinning}

You can type \cmd{@} in the question box to \textbf{pin specific papers} to a question.
A popover appears showing recently imported papers (or matching results when you keep
typing after \cmd{@}). Selecting a paper adds it as a chip above the textarea; multiple
papers can be pinned to the same question.

When at least one paper is pinned, the retrieval pipeline is \textbf{bypassed entirely}:
only the pinned papers are sent to the LLM (along with any highlights you have on them).
This is useful when you already know which paper(s) to ask about and want the answer
to be strictly grounded in those sources rather than spread across a search-ranked set.
If no paper is pinned, the normal RAG retrieval flow applies as before.

\section{Multi-Turn Chat Interface}

The Ask page has been upgraded to a \textbf{chat-style conversational interface} that supports
follow-up questions. Each turn displays the user's question and the AI's answer, with distinct
avatar indicators (blue for user, green for AI).

\litfigure{24-ask-multiturn-en.png}{Ask multi-turn conversation. User and AI messages alternate in bubbles,
each answer accompanied by source citations and a save button.}

\subsection{Conversation State Model}

The frontend maintains a \cmd{ConversationTurn[]} array, where each element contains:

\begin{itemize}
  \item \cmd{role}---\cmd{"user"} or \cmd{"assistant"}.
  \item \cmd{content}---The message text.
  \item \cmd{result}---Only present for the assistant role, containing the complete
    \cmd{AskLibraryResult} (answer, sources, model, prompt\_tokens, completion\_tokens, terms).
\end{itemize}

Conversation history is \textbf{maintained only within the current session}---it is cleared when
the page is closed or the route is changed. This is consistent with the single-turn behavior
(previous single-turn results were also not persisted across sessions).

\subsection{Multi-Turn RAG Flow}

Each follow-up question is processed using the exact same pipeline as a single turn (question
rewriting $\to$ FTS5 retrieval $\to$ evidence assembly $\to$ LLM answer). The difference lies in
how the LLM message array is constructed:

\begin{enumerate}
  \item \cmd{System prompt}---The same citation-strict system prompt as single-turn.
  \item \cmd{History messages}---Previous conversation history (all user/assistant message pairs)
    is inserted as a \cmd{ChatMessage} array between the system prompt and the current question.
  \item \cmd{Current question + evidence}---The current follow-up text plus the newly retrieved
    Sources block for this turn.
\end{enumerate}

This means:

\begin{itemize}
  \item \textbf{Context carry-over}---The LLM can see previous Q\&A, so pronouns and references
    are correctly resolved. For example, "what method did it use?"---the "it" can refer to the
    paper discussed in the previous turn.
  \item \textbf{Independent retrieval per turn}---Each follow-up re-runs question rewriting and
    FTS5 retrieval rather than reusing evidence from the first turn. This ensures that follow-ups
    can cover new semantic angles.
  \item \textbf{Independent citations}---The \cmd{[N]} numbering in each answer corresponds only
    to that turn's Sources and does not conflict with numbering from previous turns. Source cards
    are displayed independently below each answer.
\end{itemize}

\subsection{Conversation History Passing Details}

The \cmd{conversation\_history} sent to the backend is a \cmd{Vec<ConversationMessage>}, where
each element has \cmd{role} and \cmd{content}. The backend converts it into a \cmd{ChatMessage}
array and passes it to the LLM. Note: the history messages \textbf{do not contain} Sources
blocks---Sources are only attached to the current follow-up's user message. The reasons are:

\begin{itemize}
  \item Avoid token explosion---if each turn retained full Sources blocks, the context could exceed
    the model window after 5 turns.
  \item The LLM only needs the current turn's evidence to answer the current question---previous
    evidence is already reflected in previous assistant answers.
\end{itemize}

\subsection{UI Interaction}

\begin{itemize}
  \item The input box is at the bottom of the page; \kbd{Ctrl+Enter} / \kbd{Cmd+Enter} sends the
    message.
  \item The placeholder text changes to "Continue the conversation..." when conversation history
    exists.
  \item After sending a new message, the view automatically scrolls to the bottom
    (\cmd{scrollIntoView} + smooth behavior).
  \item The "Clear Conversation" button at the top resets the entire conversation history and input
    box.
  \item During loading, the AI avatar + spinning icon + "Searching..." text are displayed.
  \item Error messages appear as red text below the input box.
\end{itemize}

\begin{tipbox}
Multi-turn conversation quality depends on the context window size. If your LLM is configured with
a small context window (e.g., 4K tokens), conversations beyond 5 turns may truncate history. A
model with 8K+ window is recommended.
\end{tipbox}

\section{Save as Knowledge Note}

Each answer has a "Save as Note" button below it. Clicking generates an independent Markdown
document (at \fpath{notes/ask/<timestamp>.md}), fully preserving the \textbf{question, search
terms, answer, and evidence sources}. This note is synced along with the library during WebDAV
sync.

\begin{tipbox}
LitFolio uses a \textbf{degraded search} strategy---first rewrites the question to precise English
terms for AND matching, then falls back to raw question search, OR fuzzy matching, CJK keyword
splitting, and finally recent library papers. This means the answer \textbf{always has some papers
to work with}. If the retrieved papers are not directly relevant, the LLM honestly states they
do not address your question rather than fabricating connections. In multi-turn conversations,
follow-ups work from the same retrieved context.
\end{tipbox}

% =====================================================================
\chapter{Knowledge Graph}
\label{ch:graph}

The \cmd{/graph} route is LitFolio's literature relationship visualization hub. It aggregates
citation relationships, methodological inheritance, comparisons, and other connections between
papers---as well as associations between papers and conceptual terms---into an interactive knowledge
graph.

\section{Network Graph}

\litfigure{19-graph-network-en.png}{Knowledge Graph---network view. Blue nodes are papers, green nodes are concepts;
line colors indicate relationship types; dashed lines represent AI-discovered potential associations.}

The network graph is the default view, using force-directed layout to automatically arrange nodes:

\begin{itemize}
  \item \textbf{Paper nodes} (blue)---Fixed size; hover to display title and year.
  \item \textbf{Concept nodes} (green)---Derived from \cmd{normalized\_term} entries in the term
    table shared by two or more papers; the node label shows the number of associated papers.
  \item \textbf{Lines}---Color-coded by relationship type (see legend); solid lines are
    user-created, dashed lines are AI-discovered.
\end{itemize}

Clicking any node displays its detailed information and a list of all associated nodes in the right
sidebar.

\section{Mind Map}

\litfigure{20-graph-mindmap-en.png}{Knowledge Graph---mind map view. A concept is at the center, radiating outward;
the first ring contains associated papers; the second ring contains papers linked to those papers.}

The mind map takes a conceptual term as its root node and radiates outward in concentric rings:

\begin{itemize}
  \item \textbf{First ring}---Papers directly associated with the concept through the term table.
  \item \textbf{Second ring}---Other papers that have \cmd{paper\_links} relationships with the
    first-ring papers.
  \item Clicking the "Locate" button on a concept node switches to a mind map centered on that
    concept.
\end{itemize}

\section{Manual Relationship Creation}

The "Add Link" button on the toolbar opens a creation dialog. You need to select:

\begin{enumerate}
  \item \textbf{Source paper}---Defaults to the currently selected paper node.
  \item \textbf{Target paper}---Selected from a dropdown list.
  \item \textbf{Relationship type}---Six predefined types: \cmd{extends},
    \cmd{contradicts}, \cmd{compares}, \cmd{builds\_on}, \cmd{uses\_method}, \cmd{related}.
  \item \textbf{Description} (optional)---A brief text describing the specific connection between
    the two papers.
\end{enumerate}

User-created relationships have a fixed confidence of 1.0 and are displayed with solid lines.

\section{Citation Network}

Beyond user-created and AI-discovered relationships, LitFolio can also fetch academic citation
relationships from Semantic Scholar---showing "who the academic community considers related to this
paper." Click the "Citations" button in the paper drawer or on the Knowledge Graph page.

\litfigure{25-citations-network-en.png}{Citation network view. Papers cited by the central paper appear on the left;
papers citing it appear on the right. Papers already in the library are highlighted with a green border.}

\subsection{Data Fetching Flow}

Citation data is obtained through the Semantic Scholar Graph API:

\begin{enumerate}
  \item \textbf{Paper matching}---Send a matching request to S2 using the paper's DOI or title.
    DOI is preferred (exact match); when no DOI is available, title fuzzy matching is used, and S2
    returns the most relevant paperId.
  \item \textbf{Fetch citation relationships}---Call \cmd{GET /paper/\{id\}?fields=references,citations},
    which returns two arrays: \cmd{references} (papers this paper cites) and \cmd{citations}
    (papers that cite this paper). Each element contains title, authors, year, venue, and
    externalIds.
  \item \textbf{Caching}---Results are written to the \cmd{paper\_citations} table, with each
    record containing the direction (\cmd{references} / \cmd{citations}) and a \cmd{fetched\_at}
    timestamp. Subsequent views read directly from cache without re-calling the API.
\end{enumerate}

\subsection{Tree Display}

The citation network is displayed in a tree layout:

\begin{itemize}
  \item \textbf{Center node}---The current paper.
  \item \textbf{Left branches}---References (papers cited by this paper), up to 2 levels deep
    (the first level is direct citations; the second level is citations of those cited papers).
  \item \textbf{Right branches}---Citations (papers that cite this paper), also up to 2 levels.
  \item \textbf{Library indicator}---Papers already in your library are highlighted with a green
    border; hovering shows the TL;DR (if available).
  \item \textbf{Non-library papers}---Title, authors, and year are displayed; clicking shows
    abstract details with an "Import" button.
\end{itemize}

\subsection{Relationship with the Knowledge Graph}

The citation network and the knowledge graph are complementary dimensions:

\begin{description}
  \item[Citation Network] Objective data from the academic community---who cites whom,
    independent of your judgment.
  \item[Knowledge Graph] Derived from your subjective judgment and AI inference---methodological
    inheritance, contradictions, comparisons, and other relationships between papers.
\end{description}

The two can be cross-validated: if the knowledge graph shows that A extends B, and the citation
network confirms that A indeed cites B, the relationship's credibility is higher.

\begin{warnbox}
Semantic Scholar's citation data coverage is not 100\%---very recent papers (published less than a
week ago) or papers from niche journals may have no citation data. In such cases, the API returns
empty arrays, and the interface displays "No citation data available."
\end{warnbox}

\section{Concept Graph}

The concept graph is an advanced dimension of the knowledge graph---it visualizes not relationships
between papers, but \textbf{evolutionary relationships between methodological concepts}. Click the
"Extract Concepts" button on the toolbar to start.

\litfigure{26-concept-graph-en.png}{Concept graph. Concept nodes (green) are connected by different colored lines
indicating relationship types: replaces (pink), extends (purple), requires (orange), enables (green), competes\_with (red).}

\subsection{Concept Extraction Pipeline}

After clicking "Extract Concepts," LitFolio executes the following pipeline for each paper:

\begin{enumerate}
  \item \textbf{Text preparation}---Take the paper's TL;DR + abstract + deep-read sections (if
    available), concatenated into input text of no more than 3000 characters.
  \item \textbf{LLM extraction}---Pass the text to the LLM; the system prompt requires returning
    a JSON array where each element contains:
    \begin{itemize}
      \item \cmd{name}---Methodological concept name (e.g., \cmd{Transformer}, \cmd{LoRA},
        \cmd{Self-Attention}), with normalized naming (capitalize first letters, avoid
        abbreviations unless the abbreviation is more commonly used).
      \item \cmd{description}---A one-sentence Chinese description of the concept.
      \item \cmd{relations}---A list of relationships with other concepts, each containing
        \cmd{target} (target concept name), \cmd{relation} (relationship type), and
        \cmd{snippet} (evidence text from the original paper).
    \end{itemize}
  \item \textbf{JSON parsing}---Three-level fault tolerance (direct parsing $\to$ extract
    \cmd{[...]} substring $\to$ empty array fallback), consistent with the TL;DR parsing strategy.
  \item \textbf{Deduplication and normalization}---Convert extracted concept names to lowercase for
    deduplication comparison. If a concept with the same name already exists in the library
    (case-insensitive), reuse its existing ID; otherwise, create a new concept.
  \item \textbf{Relationship creation}---For each relationship, look up the target concept
    (name-based matching); if the target concept does not exist, create it first. Then create an
    edge in the \cmd{concept\_relations} table, recording the \cmd{evidence\_paper\_id} and
    \cmd{snippet}.
  \item \textbf{Paper-concept association}---Create a \cmd{discusses} edge in the
    \cmd{paper\_concepts} table with default relevance = 1.0.
\end{enumerate}

\subsection{Concept Relationship Types}

Five predefined relationship types, each corresponding to a different academic semantic:

\begin{longtable}{p{3cm} p{4cm} p{5.5cm}}
\toprule
\textbf{Type} & \textbf{Meaning} & \textbf{Typical Example} \\
\midrule
\endhead
\cmd{replaces} & New method replaces old method & Transformer replaces RNN \\
\cmd{extends} & Improvement on existing method & LoRA extends Fine-tuning \\
\cmd{requires} & Depends on another method & Attention requires Embedding \\
\cmd{enables} & Makes another method possible & GPU enables Deep Learning \\
\cmd{competes\_with} & Contemporary alternative & BERT competes\_with GPT \\
\bottomrule
\end{longtable}

\subsection{Graph Visualization Integration}

After extraction completes, concept nodes automatically appear in the knowledge graph:

\begin{itemize}
  \item \textbf{Network graph}---Concept nodes are green with fixed size. Relationships between
    concepts use different colored lines: replaces (pink), extends (purple), requires (orange),
    enables (green), competes\_with (red). Papers and concepts are connected by \cmd{discusses}
    edges (light green).
  \item \textbf{Mind map}---Concepts serve as root nodes radiating outward. The first ring
    contains papers directly associated with the concept (via \cmd{paper\_concepts}); the second
    ring contains papers linked to those papers (via \cmd{paper\_links}).
  \item \textbf{Legend}---The legend panel in the bottom-left corner annotates all relationship
    types with their line colors and styles.
\end{itemize}

\subsection{Manual Management}

You can manually correct AI extraction results:

\begin{itemize}
  \item \textbf{Create concept}---Manually add concepts the AI missed.
  \item \textbf{Edit concept}---Modify the name or description.
  \item \textbf{Delete concept}---Upon deletion, associated \cmd{paper\_concepts} and
    \cmd{concept\_relations} records are cascade-deleted.
  \item \textbf{Create/delete relationship}---Manually add or remove relationships between
    concepts.
  \item \textbf{Associate/disassociate papers}---Manage \cmd{discusses} edges between concepts
    and papers.
\end{itemize}

\begin{tipbox}
Concept extraction quality depends on the quality of the paper's deep-read. If a paper only has a
TL;DR without deep-read sections, the LLM has less context to work with, and the extracted concepts
may be less precise. It is recommended to run deep-read on papers before performing concept
extraction.
\end{tipbox}

\section{AI-Discovered Relationships}

Click the "AI Discover" button on the toolbar, and LitFolio will:

\begin{enumerate}
  \item Collect metadata from all papers in the library (title, year, TL;DR, methods, key findings).
  \item Group them by shared terms (up to 12 papers per batch) and send to the LLM to analyze
    potential associations.
  \item Parse the LLM's JSON output and generate relationship suggestions with confidence scores
    and descriptions for each pair of papers.
\end{enumerate}

AI-discovered relationships are displayed with dashed lines and have confidence scores below 1.0.
In the reader's term panel, AI relationships show a confidence percentage. You can:

\begin{itemize}
  \item \textbf{Accept}---Promote the AI relationship to a user relationship (confidence becomes
    1.0).
  \item \textbf{Reject}---Delete the relationship suggestion.
\end{itemize}

\section{Filtering and Legend}

The toolbar provides two filter groups:

\begin{itemize}
  \item \textbf{Relationship type}---Filter which lines are displayed by type label.
  \item \textbf{Include concepts}---Toggle concept node visibility. When disabled, only direct
    paper-to-paper relationships are shown.
\end{itemize}

The legend panel in the bottom-left corner annotates node colors and the line style for each
relationship type.

\section{Reader Integration}

At the bottom of the reader's term panel (Chapter~\ref{ch:reader}), all associated papers for the
current paper are displayed. Each association shows the relationship type and source (user/AI);
clicking navigates to the corresponding paper's reader view.

\begin{tipbox}
The knowledge graph data comes entirely from the local SQLite database (intersection of the
\cmd{paper\_links} table and the \cmd{paper\_terms} table). Relationship data is synced along
with the library during WebDAV sync.
\end{tipbox}

% =====================================================================
\chapter{Similar Paper Recommendations}
\label{ch:similar}

After reading a paper, the natural question is "what else is worth reading?" Similar Paper
Recommendations use Semantic Scholar's recommendation API to automatically find related literature.

\section{Usage}

Click the "Similar Papers" button in the paper drawer or at the top of the reader. LitFolio uses
the paper's DOI or title to send a recommendation request to Semantic Scholar.

\litfigure{28-similar-papers-en.png}{Similar Paper Recommendations panel. Result cards display title, authors, year,
abstract snippet, and relevance score.}

\subsection{Recommendation Fetching Flow}

\begin{enumerate}
  \item \textbf{Paper matching}---Use DOI (preferred) or title to match a paperId with Semantic
    Scholar. The same matching logic as the citation network is used.
  \item \textbf{Call recommendation API}---Request
    \cmd{GET /paper/\{id\}/recommendations}, which returns a list of papers sorted by relevance.
    Each entry contains title, authors, year, abstract, and citationCount.
  \item \textbf{Library filtering}---Iterate through the recommendation results and check each
    against the library using DOI and arXiv ID. Papers already in the library are removed from the
    results.
  \item \textbf{Caching}---Filtered results are written to the \cmd{recommendation\_cache} table
    with a \cmd{fetched\_at} timestamp. Results are read directly from cache for 7 days.
\end{enumerate}

\subsection{Result Display}

Each result card contains:

\begin{itemize}
  \item Title, authors (first 3 + \cmd{et al.}), year.
  \item Abstract snippet (first 200 characters).
  \item Citation count (if S2 returned this field).
  \item "Import" button---clicking fetches full metadata and attempts to download the PDF.
\end{itemize}

\subsection{Edge Cases}

\begin{itemize}
  \item \textbf{Papers without DOI}---Title fuzzy matching may return an incorrect paperId,
    leading to irrelevant recommendation results. In this case, the interface displays "Match
    accuracy is low; results are for reference only."
  \item \textbf{Very recent papers}---S2's recommendation model requires papers to have some
    citation accumulation before generating recommendations. Papers published less than a month ago
    typically have no recommendation results.
  \item \textbf{Niche fields}---S2's data coverage is not 100\%; some non-English or
    non-mainstream conference papers may not match.
\end{itemize}

\begin{tipbox}
If recommendation results are empty, try using the manual search in "Topic Discovery" (Chapter 7),
or check AI-discovered associations in the Knowledge Graph.
\end{tipbox}

% =====================================================================
\chapter{Literature Review Draft Generation}
\label{ch:litreview}

After collecting 20--50 papers on a research direction, the first draft of a literature review is
the hardest to write. LitFolio's AI review generation helps you build the skeleton.

\section{Usage}

Select multiple papers in the library (using checkbox mode) and click "Generate Review" on the
toolbar. You can also select "Generate Review" from a folder's right-click context menu.

\litfigure{29-lit-review-en.png}{Generated literature review draft. Organized by theme sections, each citing relevant
papers, with a research gaps section at the end.}

\subsection{Generation Pipeline}

Review generation is a two-step pipeline:

\subsubsection{Step 1: Grouping}

Send the selected papers to the LLM and ask it to group them by theme/method/application area.
The LLM returns JSON:

\begin{verbatim}
{
  "groups": [
    {
      "name": "Attention Mechanism Variants",
      "paper_ids": ["p1", "p3", "p7"],
      "rationale": "These papers all improve attention computation efficiency"
    },
    {
      "name": "Positional Encoding Methods",
      "paper_ids": ["p2", "p5"],
      "rationale": "Addressing Transformer position awareness"
    }
  ]
}
\end{verbatim}

Grouping strategy options: \cmd{method} (by method, default), \cmd{year} (by era),
\cmd{application} (by application area). Different strategies produce different groupings.

\subsubsection{Step 2: Writing}

Concatenate the grouping results with each paper's TL;DR + deep-read sections + abstract into a
prompt, and ask the LLM to write the review. System prompt constraints:

\begin{itemize}
  \item Each group corresponds to a second-level section.
  \item Every paper is cited at least once in the text, using the format \cmd{[Author, Year]}.
  \item The introduction outlines the research background and the review's scope.
  \item A "Research Gaps" section at the end identifies unresolved problems in the current
    literature.
  \item The conclusion summarizes main findings and future directions.
  \item Language is academic Chinese, consistent with the style of formal review papers.
  \item Do not fabricate papers or data not provided.
\end{itemize}

\subsection{Input Materials}

Materials provided to the LLM for each paper (by priority, total tokens not exceeding the window
limit):

\begin{enumerate}
  \item TL;DR + key findings (approx.\ 100 tokens/paper)
  \item Deep-read sections: Problem / Method / Comparison / Limitations (approx.\ 400 tokens/paper)
  \item Abstract (approx.\ 200 tokens/paper)
\end{enumerate}

Input for 20 papers is approximately 14,000 tokens; including the system prompt and output space,
a model with at least a 16K context window is required.

\section{Output and Editing}

\begin{itemize}
  \item Output is in Markdown format, 2000--5000 words (for 20 papers).
  \item Can be edited directly in the interface to correct AI summaries or supplement your own
    observations.
  \item Can be saved as a note file (\fpath{notes/review/<timestamp>.md}).
  \item Can be regenerated---choose a different grouping strategy (by method / by year / by
    application area).
\end{itemize}

\begin{warnbox}
The review draft is a starting point, not a finished product. The AI may overlook subtle
differences between papers or overgeneralize. Be sure to manually review and supplement with your
own analysis. Pay particular attention to whether the "Research Gaps" section covers the directions
you consider important---the AI may over-emphasize certain papers while neglecting others.
\end{warnbox}
